\begin{theorem}[G-21: Maximal permanent-image dimension in every larger ambient space]
\label{G-21}
Let $m\ge17$ and $r\ge4$. Let $\mathcal P_{m,r}$ be the
Zariski closure of the forms obtained as permanents of
$m\times m$ pencils of linear forms in $r$ variables over
$\mathbb C$. Then
\begin{equation}\label{G21:dimension}
 \dim\mathcal P_{m,r}=rm^2-2m+2.
\end{equation}
Thus four variables already suffice to attain the full
parameter count modulo the row--column scaling directions;
each additional variable contributes exactly $m^2$ dimensions.
For a general coefficient specialization
$R=\mathbb C[X_{ij}]\longrightarrow S_r=\mathbb C[x_1,\ldots,x_r]$,
the corresponding graded first homology satisfies
\begin{equation}\label{G21:Tor}
 \operatorname{Tor}_1^R(R/J_m,S_r)_d=0
 \qquad(d\le m),
\end{equation}
where $J_m$ is the submaximal-permanent ideal used in R-23.
The assertion in \eqref{G21:Tor} concerns the specified graded
pieces and does not assert exactness of the entire restricted
resolution in higher degrees.
\end{theorem}

\begin{proof}
For a quaternary pencil $A$, write
$G_{ij}(A)=\partial\operatorname{per}(A)/\partial A_{ij}$.
R-23's cyclic-pencil proof shows that the $m^2$ degree-$(m-1)$
forms $G_{ij}(A)$ are linearly independent on a nonempty
Zariski open subset of the quaternary parameter space. R-27
shows that the differential image
\[
 \operatorname{span}\{x_kG_{ij}(A):1\le k\le4,\ 1\le i,j\le m\}
\]
has dimension $4m^2-2m+2$ on another nonempty open subset.
Their parameter space is irreducible, so choose one $A$
satisfying both conditions.

Regard this same pencil as an $r$-variable pencil with zero
coefficients on $x_5,\ldots,x_r$. The differential in every
new coefficient direction has the form $x_kG_{ij}(A)$,
$5\le k\le r$. For each fixed new variable these $m^2$ forms
are independent, because multiplication by $x_k$ is injective.
Their spans for different new variables, and the original
four-variable differential span, have disjoint multidegrees
in $x_5,\ldots,x_r$. Consequently the rank at this explicit
kind of parameter point is
\[
 (4m^2-2m+2)+(r-4)m^2=rm^2-2m+2.
\]
The generic image dimension is at least this rank. The
row--column torus has the same effective dimension $2m-2$
in every variable count, and gives the opposite upper bound
$rm^2-2m+2$, exactly as in R-13/R-23. This proves
\eqref{G21:dimension}. The use of a pencil with initially
unused variables is legitimate: a nonzero coefficient minor
at that point defines a nonempty open subset of the full
$r$-variable parameter space.

For the Tor statement, use the same first part of the graded
resolution and the same cited linear-syzygy theorem as in
R-23. In degree $m$, the relevant complex after specialization is
\[
 \mathbb C^{2m-2}\longrightarrow (S_r)_1^{m^2}
       \longrightarrow (S_r)_m.
\]
The old row--column relations remain independent on a nonempty
open subset: they are independent at the four-variable pencil
just used, regarded as an $r$-variable pencil. The last map
is the pencil differential and has the exact generic rank
\eqref{G21:dimension}. Its kernel therefore has dimension
$rm^2-(rm^2-2m+2)=2m-2$ and equals the image of the first
map. This proves the degree-$m$ assertion.
In degree $m-1$, the last map sends constants to the independent
$m^2$ cofactors. Below that degree its source is zero. These
give all lower-degree assertions in \eqref{G21:Tor}.
This tensor-complex argument remains meaningful when $r>m^2$;
in that case the coefficient map is an algebra homomorphism
into a larger polynomial ring, and is not called a quotient
by $m^2-r$ linear equations.
\end{proof}

\paragraph{Exact scope.}
This gives an image-dimension and low-degree derived-restriction
statement in all $r\ge4$. It does not compute individual
$r$-row multiplicities, identify all higher syzygies, or
turn the dimension statement into a polynomial-time positive
formula in growing $r$. Those are distinct remaining tasks.
